%% notes-tex/microbe-perturb-seq/sections/1-introduction.tex
%% SOURCE: hand-authored prose. No generating script, and there should not be one.
%% Numbers quoted here are cross-referenced to the sections that derive them.

\section{Introduction\secstatus{todo}}
\label{sec:intro}

Metabolic engineering builds strains by stacking perturbations, but the design
decisions are made from data on single perturbations measured one at a time.
Pooled CRISPR interference screens read out by single-cell RNA sequencing
(\term{Perturb-seq}) close that gap: a library of guide RNAs (\term{gRNA}s) is
transformed into a population, every cell is sequenced, and each cell's gRNA
identity is recovered from its own transcriptome, so one experiment yields a
matrix from perturbation to full transcriptional response rather than a list of
fitness hits~\citep{dixitPerturbSeqDissectingMolecular2016}. The genome-wide
MAGIC library this design would screen returns exactly such a list, scored by
growth~\citep{lianMultifunctionalGenomewideCRISPR2019}; the same library read out
by Perturb-seq returns the matrix instead. That matrix is the highest-density
source of genotype-to-expression training data available.

Two properties of the readout are worth separating, because running them
together is what makes the advantage sound larger than it is. The
\emph{transcriptional} readout is what scores perturbations a growth screen
cannot see: a gene can be fitness-neutral and still move hundreds of transcripts.
\emph{Single-cell} resolution is what makes a pooled library readable at all,
because it assigns each transcriptome to the gRNA that produced it. The effect
estimate is not per cell: cells sharing a gRNA are summed into one pseudobulk
profile, so precision is set by how many cells share a gRNA rather than by how
deeply any one of them was read (\cref{sec:cells-per-pert}).

Three developments have made the yeast version tractable. A split-pool barcoding
chemistry removes the instrument from the cell-isolation
step~\citep{brettnerUltraHighthroughputMassively2024}; pooled CRISPR activation
and interference have been joined to single-cell transcriptomics in bacteria,
the first single-cell CRISPR screen in that
domain~\citep{brandnerPooledSinglecellCRISPRa2025}; and a genome-scale yeast
deletion collection has been rebuilt so its genotype tag is transcribed and
therefore visible to the
assay~\citep{nadal-ribellesSinglecellResolvedGenotypephenotype2025}. What follows
establishes what such a screen costs in \org{S.\ cerevisiae}, how many genes can
be perturbed simultaneously before the measurement stops being interpretable, and
which parameter is actually binding.

Two further questions decide what gets run first, and a cost table answers
neither. \Cref{sec:scaling} asks how the design grows along the two axes a
metabolic engineering campaign cares about, more perturbations per cell and more
environments per experiment. Cell demand grows quadratically in the first and
linearly in the second, which is an argument for spending on environments rather
than a verdict: a gene combination and a growth condition are different kinds of
information, and a model that proposes edits needs both.
\Cref{sec:host} asks what changes when the host is \org{P.\ kudriavzevii}. Almost
none of the genetics assumed here is published for that organism, which is not
the same as unavailable: an assembled genome and a working CRISPRi system exist
in house and are not in the
literature\external{in-house, unpublished; genome from JGI}.

The first experiment therefore carries no perturbations at all, and that is a
deliberate choice rather than a forced one. Nobody on this project has run the
technology before, so an unperturbed run buys footing: it answers the three
things that are genuinely uncertain about moving host, wall digestion, UMIs
recovered per cell and ribosomal fraction, while depending on no genetics at all
(\cref{sec:host}). If one unperturbed run looks wasteful for the cells it
consumes, the way to spend it better is across several environments, not by
adding guides to it.
